\subsection{Knowledge Graphs}
Knowledge Graphs (KGs) are semantic data structures that represent information through entities and the relationships between them. Unlike traditional relational databases, Knowledge Graphs organize data as a graph, where nodes represent entities and edges represent the semantic relationships connecting them. This representation enables the integration of heterogeneous data sources while preserving their semantic meaning.

Knowledge Graphs have become increasingly important in domains such as search engines, healthcare, recommendation systems, and artificial intelligence because they facilitate semantic querying, knowledge discovery, and reasoning. In this thesis, Knowledge Graphs provide the conceptual foundation for constructing a Virtual Knowledge Graph over a relational movie database, enabling users to access data through semantic queries rather than directly interacting with SQL tables.
\subsection{Virtual Knowledge Graphs}
Virtual Knowledge Graphs (VKGs) extend the concept of Knowledge Graphs by providing a semantic layer over existing relational databases without physically converting the data into RDF. Instead of storing the information in a graph database, the data remain in the original relational database while the ontology and semantic mappings create a virtual RDF representation.

This approach is based on the Ontology-Based Data Access (OBDA) paradigm, where users interact with the data through the ontology rather than directly accessing the relational database. When a SPARQL query is submitted, it is automatically translated into an SQL query and executed on the underlying database. Therefore, the RDF graph is generated virtually at query time, avoiding data duplication while preserving the efficiency of relational database systems.

In this thesis, the Virtual Knowledge Graph is constructed over a PostgreSQL database containing movie-related data. The ontology provides the conceptual representation of the domain, while Ontop uses semantic mappings to connect ontology terms with the corresponding relational tables. As a result, users can retrieve information using SPARQL without directly interacting with the relational database.

\subsection{Ontology}
An ontology provides a formal and explicit specification of the concepts within a domain and the relationships between them. In Semantic Web technologies, ontologies define a shared vocabulary that enables both humans and machines to interpret data consistently. An ontology typically consists of classes, object properties, data properties, and individuals, together with logical constraints such as domains and ranges.

Unlike relational databases, which focus on storing data efficiently, ontologies focus on representing the meaning of data. They provide a conceptual view of a domain independently of how the information is physically stored. Consequently, the ontology does not contain the actual data; instead, it describes the structure and semantics of the information.

In this thesis, the ontology models the movie domain by representing entities such as movies, users, reviews, and watch history together with the semantic relationships between them. The ontology serves as the conceptual layer of the Virtual Knowledge Graph and enables semantic querying through SPARQL.

\subsection{Resource Description Framework (RDF)}

The Resource Description Framework (RDF) is the standard data model used to represent semantic information on the Semantic Web. RDF expresses knowledge as a collection of triples, where each triple consists of a subject, a predicate, and an object. This simple structure makes it possible to describe entities and their relationships in a machine-readable and interoperable format.

In an RDF graph, subjects and objects are represented as nodes, while predicates define the relationships between them. RDF therefore provides the underlying representation used by Knowledge Graphs and ontologies.

In this thesis, the ontology and the semantic mappings together produce a virtual RDF graph over the relational PostgreSQL database. Although the data remain physically stored in relational tables, they are exposed as RDF triples through Ontop, enabling semantic querying using SPARQL.

\subsection{SPARQL}
SPARQL (SPARQL Protocol and RDF Query Language) is the standard query language for retrieving and manipulating data represented in RDF. Similar to SQL for relational databases, SPARQL enables users to query semantic data by specifying graph patterns rather than relational tables.

A SPARQL query matches patterns of RDF triples and returns the corresponding results from the Knowledge Graph. Depending on the query, the results may include entities, relationships, or specific attribute values.

In an Ontology-Based Data Access (OBDA) system, users formulate queries in SPARQL without needing to know the underlying relational database schema. Systems such as Ontop automatically translate SPARQL queries into SQL queries, execute them on the relational database, and return the results as if they had been retrieved from an RDF graph.

In this thesis, SPARQL is used to validate the Virtual Knowledge Graph and to retrieve information from the PostgreSQL database through the ontology and semantic mappings.

\subsection{Ontology-Based Data Access (OBDA)}
Ontology-Based Data Access (OBDA) is an approach that enables users to access relational data through a semantic layer represented by an ontology. Instead of interacting directly with database tables and writing SQL queries, users formulate their queries using the concepts and relationships defined in the ontology.

The ontology provides a conceptual view of the underlying data, while semantic mappings establish the correspondence between ontology elements and the relational database. When a user submits a SPARQL query, the OBDA system automatically translates it into an equivalent SQL query, executes it on the relational database, and returns the results through the semantic layer.

One of the main advantages of OBDA is that it allows users to benefit from semantic querying without physically transforming relational data into RDF. This approach preserves the efficiency and scalability of relational database systems while providing a more intuitive and meaningful way to access information.

In this thesis, the proposed framework adopts the OBDA paradigm to provide semantic access to a PostgreSQL movie database. The ontology, mappings, and Ontop platform together enable the construction of a Virtual Knowledge Graph without duplicating the original data.
\subsection{Ontop}
Ontop is one of the most widely used platforms for implementing Ontology-Based Data Access (OBDA). It enables users to access relational databases through an ontology without physically transforming the data into RDF. Instead, Ontop creates a Virtual Knowledge Graph by connecting the ontology to the relational database through semantic mappings.

One of the key features of Ontop is its ability to automatically translate SPARQL queries into SQL queries. When a user submits a SPARQL query, Ontop uses the semantic mappings to identify the corresponding database tables and attributes, generates an equivalent SQL query, executes it on the relational database, and returns the results as RDF.

This approach combines the expressive power of semantic technologies with the efficiency of relational database systems. Since the data remain stored in the original database, no data duplication is required, and the Virtual Knowledge Graph is generated dynamically at query time.

In this thesis, Ontop is used to connect the PostgreSQL movie database with the ontology through semantic mappings, allowing users to retrieve information using SPARQL while the data remain stored in the relational database.
\subsection{Large Language Models (LLMs)}
Large Language Models (LLMs) are artificial intelligence models trained on vast amounts of textual data to understand and generate natural language. Based on transformer architectures, LLMs have demonstrated remarkable capabilities in tasks such as question answering, text generation, summarization, reasoning, and code generation.

Recently, LLMs have also been applied to Semantic Web and Knowledge Graph applications. They can assist users by interpreting natural language questions, generating SPARQL queries, supporting ontology engineering, and facilitating interaction with semantic data without requiring expertise in formal query languages.

In this thesis, LLMs are investigated as a component of the proposed framework to improve user interaction with the Virtual Knowledge Graph. In particular, they are used to interpret user requests, support intent detection, and enable the integration of semantic querying with recommendation functionalities. The design and implementation of the proposed LLM-based framework are presented in Chapter~4.
This chapter introduced the theoretical foundations required for the remainder of the thesis. It presented the fundamental concepts of Knowledge Graphs, Virtual Knowledge Graphs, ontologies, RDF, SPARQL, Ontology-Based Data Access, Ontop, and Large Language Models. Together, these concepts establish the theoretical foundation for the framework proposed in this thesis, which is presented in the following chapters.
